
%\begin{abstract}
Counter machines are a particularly simple, Turing-complete
model of computation.
We study the expressive power of nondeterministic
counter machines which are symmetric, namely their
transitions are invariant under swapping of counters.
As the main result we demonstrate that symmetric 2-counter machines
are not Turing-complete, by proving decidability of their 
reachability problem.
%Our approach builds on specific properties of 2-counter
%machines, and does not extend to 3 or more counters.
We conjecture the decidability 
for nondeterministic machines with arbitrarily many counters,
as long as (their transitions) are invariant under all permutations
of counters. 
We leave the status for symmetric 3-counter machines open, but
identify a bound on the impact of symmetry by proving
undecidabilty for 3-counter machines invariant under the
cyclic permutation group.
%\end{abstract}

\bigskip

\section{Introduction}

Counter machines (\cm{s}) are a particularly simple
model of computation \cite{Minsky67}.
A counter machine has a number of registers where it
stores its data, each register storing a natural number.
While the particular instruction sets may vary, 
common machine instructions are: counter increment, 
counter decrement, and counter zero-test. 
Many variants of the model have been introduced and studied,
and both deterministic and nondetermistic 
machines are widely studied in the literature 
(often under different names, for instance 
\emph{counter automata}).
Numerous syntactic restrictions of the model are also widely
studied \cite{HKOW09,Ibarra78,HopcroftPansiot,FS00}
since,
as shown already by Minsky in the early 1960s \cite{Minsky61},
in its full generality the model is Turing-complete:
deterministic 2-counter machines can faithfully simulate Tuting machines
and have undecidable halting problem. 

The purpose of this paper is to understand the impact of
\emph{symmetry} on the expressive power of \cm{s}. %counter machines.
We investigate the \emph{nondeterministic} model, where
a machine may choose between different istructions executable
in the same state.

\para{Symmetric counter machines}
As the starting point, we specifically choose the set of 
machine instructions:
counter \emph{updates} that simultaneously modify all counters  
by given numbers, and counter \emph{zero-tests} that are 
executable only if a given counter has value zero.
Every \emph{transition} $(q,i,q')$ of a machine changes its state
from $q$ to $q'$,
and executes an attached instruction $i$.
Executing a transition in a given configuration constitutes a 
\emph{step} of a machine.
The model we investigate is thus essentially \vass
(vector addition systems with states) \cite{HopcroftPansiot}
extended with zero-tests, and
is clearly equivalent to standard \cm{s}. %counter machines.
In this paper we mostly focus on \cm{s} %counter machines 
that are 
\emph{symmetric}, by which we mean that the set of transitions
is invariant under all permutations $\sigma$ of counters:
for every transition $(q,i,q')$, the machine has also a transition
$(q,\sigma(i), q')$, where $\sigma(i)$ is the result of 
applying $\sigma$ to the counters in the instruction $i$.
Specifically, we choose to impose the invariance only on update transitions
(when $i$ is a counter update), but
not on zero-test transitions, since the latter does not seems
necessary for our proofs to work.
Here is an illustrative example of a 2-\cm (a \cm having two counters)
with
two control states and
six transitions, including five symmetric update transitions
and one zero-test transition:
\sla{the example to be improved}
% \begin{figure}
\begin{center}
\scalebox{0.7}
{\begin{tikzpicture} [draw=black!70,
  fill=blue!20, node distance = 4cm, on grid, auto,
  every initial by arrow/.style = {thin}]
% State q0 
\node (q0) [state, scale=1.4, fill=blue!20, initial text = {}] {$p$};
% State q1  
\node (q1) [state, scale=1.4, fill=blue!20, right = of q0] {$q$};
% Arrows
\path [-stealth, thick]
   (q0) edge[bend left] node[scale=1.5] {\tiny{$\textbf{zero-test}_1$}}  (q1)
   (q1) edge[bend left] node[scale=1.1] {$(-1,-1)$}  (q0)
   (q0) edge [loop above] node[scale=1.5] {$\substack{(-1, 3) \\ (3, -1)}$}()
   (q1) edge [loop above] node[scale=1.5] {$\substack{(1, -1) \\ (-1, 1)}$}();
\end{tikzpicture}}
\end{center}
% \caption{A symmetric 2-counter machine.}
\label{fig:example}
% \end{figure}


\para{Contribution}
The initial motivation of this study is the immense impact of symmetry
on complexity of the \vass reachability problem:
while the problem is {\sc Ackermann}-complete in general \cite{LS19,CO21,L21},
the complexity drops to
only \pspace-complete for symmetric \vass \cite{KL25}.
As our main result, we confirm the immense drop in case of 
2-\cm{s}
(in this case, symmetricity amounts to invariance under all swaps of 
counters).
We prove decidability of the configuration-to-configuration reachability
problem, when one aks if a given machine has a sequence of steps
starting in a given source configuration and ending in
a given target one:

\begin{theorem} \label{thm:main}
The reachability problem  is decidable for symmetric 2-counter machines.
\end{theorem}

Symmetric two-counter machines are thus not Turing-complete.
We find this result quite surprising, as it applies to a model
with unrestricted zero-test instructions.

\begin{remark}
In the proof we do not need symmetry of zero-test transitions,
only symmetry of updating transitions.
Interestingly, the converse restriction, namely imposing symmetry on
zero-test transitions but not on update transitions immediately 
leads to undecidability.
\end{remark}

A natural question is whether \Cref{thm:main} extends to larger numbers of
counters. 
In particular, is the reachability problem still decidable
for symmetric 3-counter machines, whose update transitions are invariant
under all 6 permutations of the 3 counters?
Our decidability argument does not extend to the 3-counter setting, 
as the argument relies on a
number of low-dimensional properties of reachable sets and reachability
relations.
We thus leave this intriguing question for further research, conjecturing
decidability for any numner of counters:

\begin{conjecture}
The reachability problem  is decidable for symmetric counter machines.
\end{conjecture}

Another question we leave open is the exact complexity of the problem.
As the lower bound, we known nothing more than 
\pspace-hardness inherited from symmetric 2-\vass \cite{KL24}.
It is possible that the reachability problem is \pspace-complete
in symmetric $d$-\cm{s} for every $d\leq 2$.
\sla{expand}

\medskip

It is natural to consider invariance under certain subgroups
of permutations: for every subgroup $G\leq \symg d$
one can consider the class of \emph{$G$-counter machines}
($G$-\cm{s}), namely 
$d$-counter machines invariant under all permutations from $G$.
On one end, 
when $G=\symg d$, the general definition instantiates to 
symmetric $d$-\cm{s}.
On the opposite end, when $G$ is the trivial group one gets classical counter
machines.
Between the two extremities there is a landscape of models, 
whose expressive power is related according to the permutation subgroup relation:
whenever $G\leq H$, every $H$-\cm is automatically a
$G$-\cm.
For instance, every $\symg 3$-\cm is automatically
a $\cycg 3$-\cm, for the cyclic 
permutation group $\cycg 3 \leq \symg 3$.
As an additional result, we identify a bound on the impact of symmetry,
already for 3-counter machines, namely we prove undecidability 
for $\cycg 3$-counter machines:

\begin{theorem} \label{thm:undecid}
The reachability problem  is undecidable for $\cycg 3$-counter machines.
\end{theorem}


\para{Related research}

Identifying decidable restrictions of counter machines
is a very active line of research since 1960s, for instance 
1-counter, reversal-bounded, or restricted 2-counter machines 
\cite{HKOW09,Ibarra78,FS00}.
A prominent example is counter machines without zero-tests, that is
vector addition systems with states
(\vass), shown to have decidable reachability in early 1980s
\cite{Mayr81,Kosaraju82,Lambert92}.
The decidability extends to 
\cm{s}
with only one zero-tested counter \cite{Bonnet11},
or with with nested zero-tests \cite{Reinhardt08,GCL25}.
On the other hand, most of natural restrictions of the counter model
preserves Turing-completeness, as long as no constraint is imposed
on zero-tests.
For instance, considering deterministic \cm{s}, 
the reachability problem is still undecidable
in the subclass of \emph{reversible} \cm{s},
where the step relation is reverse deterministic, and hence every step
may be uniquely undone \cite{Morita96}.

In the special case of 1 counter, symmetric 1-\cm{s} coincide with 
general 1-\cm{s}.
Their reachability problem is \np-complete 
if transition updates are encoded in binary, and \nl-complete 
in unary encoding \cite{HKOW09}.

% (Symmetric) 2-counter machines are an extesion of (symmetric) 2-\vass by 
% zero-tests.
Restricting the use of zero-tests in 2-counter machines often
leads not only to decidability, but even 
to \pspace-completenss of the reachability problem.
Assuming binary encoding,
the \pspace-completeness has been obtained for 2-\vass 
\cite{Blondin15,BEFGHLMT21}
(assuming unary encoding, we have \nl-completeness \cite{ELT16}),
but also for the further restriction to symmetric 2-\vass \cite{KL24},
and for a relaxation of 2-\vass to 
2-counter machines with just one zero-tested counter \cite{FLS18}.
In all the three cases the reachability relation is effectively semilinear, which
is a key for obtaining the \pspace upper complexity bound, and which 
fails when the number of counters is at least 3 (even for 3-\vass) 
\cite{HopcroftPansiot}.
We do not know the exact complexity of symmetric 2-counter machines,
but it is possible that the complexity remains \pspace.
Furthermore, we do not know if the reachability relation of a
symmetric 2-counter machine is semilinear, but we believe it is so.

Contrarily to the avove-mentioned restrictions of zero-tests, 
our model leaves the zero-tests unrestricted, which makes
\Cref{thm:main} particularly interesting.
We conclude by recalling another recently obtained decidability result that
also does not assume any explicit restriction of zero-tests \cite{Dudenhefner22}.
% , and hence seems as surprising as our \Cref{thm:main}.
It amounts to observing that
the above-mentioned Turing-completeness of 
reversible counter machines seems to be very fragile with respect
to the choice of instruction set of a \cm.
As shown in \cite{Dudenhefner22},
for certain 
choice of instruction set, the Turing-complete model of 2-\cm{s}
becomes decidable under the reversibility restriction.
% 
